%-----------------------------------------------------------------------------------------
% Autor dieser Vorlage:
% Stefan Macke (http://fachinformatiker-anwendungsentwicklung.net)
% Permalink zur Vorlage: http://fiae.link/LaTeXVorlageFIAE
%
% Sämtliche verwendeten Abbildungen, Tabellen und Listings stammen von Dirk Grashorn.
%
% Lizenz: Creative Commons 4.0 Namensnennung - Weitergabe unter gleichen Bedingungen
% -----------------------------------------------------------------------------------------

\documentclass[
	12pt,
	ngerman,
	toc=listof, % Abbildungsverzeichnis sowie Tabellenverzeichnis in das Inhaltsverzeichnis aufnehmen
	toc=bibliography, % Literaturverzeichnis in das Inhaltsverzeichnis aufnehmen
	footnotes=multiple, % Trennen von direkt aufeinander folgenden Fußnoten
	parskip=half, % vertikalen Abstand zwischen Absätzen verwenden anstatt horizontale Einrückung von Folgeabsätzen
	numbers=noendperiod % Den letzten Punkt nach einer Nummerierung entfernen (nach DIN 5008)
]{scrartcl}
\pdfminorversion=5 % erlaubt das Einfügen von pdf-Dateien bis Version 1.7, ohne eine Fehlermeldung zu werfen (keine Garantie für fehlerfreies Einbetten!)
\usepackage[utf8]{inputenc} % muss als erstes eingebunden werden, da Meta/Packages ggfs. Sonderzeichen enthalten

% !TEX root = Projektdokumentation.tex

% Hinweis: der Titel muss zum Inhalt des Projekts passen und den zentralen Inhalt des Projekts deutlich herausstellen
\newcommand{\titel}{Entwicklung von SSRS Statistics}
\newcommand{\untertitel}{Webanwendung zur Visualisierung von Metadaten zu den SSRS}
\newcommand{\kompletterTitel}{\titel{} -- \untertitel}

\newcommand{\autorName}{Yannik Queisler}
\newcommand{\autorAnschrift}{Josef-Paris-Straße 17}
\newcommand{\autorOrt}{45721 Haltern am See}

\newcommand{\betriebLogo}{LogoBetrieb.pdf}
\newcommand{\betriebName}{Kassenärztliche Vereinigung Westfalen-Lippe}
\newcommand{\betriebAnschrift}{Robert-Schimrigk-Straße 4-6}
\newcommand{\betriebOrt}{44141 Dortmund}

\newcommand{\ausbildungsberuf}{Fachinformatiker für Anwendungsentwicklung}
\newcommand{\betreff}{Dokumentation zur betrieblichen Projektarbeit}
\newcommand{\pruefungstermin}{Sommer 2019}
\newcommand{\abgabeOrt}{Dortmund}
\newcommand{\abgabeTermin}{20.05.2019}
 % Metadaten zu diesem Dokument (Autor usw.)
% !TEX root = ../Projektdokumentation.tex

% Anpassung an Landessprache ---------------------------------------------------
\usepackage{babel}

% Umlaute ----------------------------------------------------------------------
%   Umlaute/Sonderzeichen wie äüöß direkt im Quelltext verwenden (CodePage).
%   Erlaubt automatische Trennung von Worten mit Umlauten.
% ------------------------------------------------------------------------------
\usepackage[T1]{fontenc}
\usepackage{textcomp} % Euro-Zeichen etc.

% Schrift ----------------------------------------------------------------------
\usepackage{lmodern} % bessere Fonts
\usepackage{relsize} % Schriftgröße relativ festlegen

% Zitate -----------------------------------------------------------------------
\usepackage{csquotes}
\usepackage[begintext=„, endtext=”]{quoting}
\SetBlockEnvironment{quoting} % Specify quoting

% Tabellen ---------------------------------------------------------------------
\PassOptionsToPackage{table}{xcolor}
\usepackage{tabularx}
% für lange Tabellen
\usepackage{longtable}
\usepackage{array}
\usepackage{ragged2e}
\usepackage{lscape}
\newcolumntype{w}[1]{>{\raggedleft\hspace{0pt}}p{#1}} % Spaltendefinition rechtsbündig mit definierter Breite

% Grafiken ---------------------------------------------------------------------
\usepackage[dvips,final]{graphicx} % Einbinden von JPG-Grafiken ermöglichen
\usepackage{graphics} % keepaspectratio
\usepackage{floatflt} % zum Umfließen von Bildern
\graphicspath{{Bilder/}} % hier liegen die Bilder des Dokuments

% Sonstiges --------------------------------------------------------------------
\usepackage[titles]{tocloft} % Inhaltsverzeichnis DIN 5008 gerecht einrücken
\usepackage{amsmath,amsfonts} % Befehle aus AMSTeX für mathematische Symbole
\usepackage{enumitem} % anpassbare Enumerates/Itemizes
\usepackage{xspace} % sorgt dafür, dass Leerzeichen hinter parameterlosen Makros nicht als Makroendezeichen interpretiert werden

\usepackage{makeidx} % für Index-Ausgabe mit \printindex
\usepackage[printonlyused]{acronym} % es werden nur benutzte Definitionen aufgelistet

% Einfache Definition der Zeilenabstände und Seitenränder etc.
\usepackage{setspace}
\usepackage{geometry}

% Symbolverzeichnis
\usepackage[intoc]{nomencl}
\let\abbrev\nomenclature
\renewcommand{\nomname}{Abkürzungsverzeichnis}
\setlength{\nomlabelwidth}{.25\hsize}
\renewcommand{\nomlabel}[1]{#1 \dotfill}
\setlength{\nomitemsep}{-\parsep}

\usepackage{varioref} % Elegantere Verweise. „auf der nächsten Seite“
\usepackage{url} % URL verlinken, lange URLs umbrechen etc.

\usepackage{chngcntr} % fortlaufendes Durchnummerieren der Fußnoten
% \usepackage[perpage]{footmisc} % Alternative: Nummerierung der Fußnoten auf jeder Seite neu

\usepackage{ifthen} % bei der Definition eigener Befehle benötigt
\usepackage{todonotes} % definiert u.a. die Befehle \todo und \listoftodos
\usepackage[square]{natbib} % wichtig für korrekte Zitierweise

% PDF-Optionen -----------------------------------------------------------------
\usepackage{pdfpages}
\pdfminorversion=5 % erlaubt das Einfügen von pdf-Dateien bis Version 1.7, ohne eine Fehlermeldung zu werfen (keine Garantie für fehlerfreies Einbetten!)
\usepackage[
    bookmarks,
    bookmarksnumbered,
    bookmarksopen=true,
    bookmarksopenlevel=1,
    colorlinks=true,
% diese Farbdefinitionen zeichnen Links im PDF farblich aus
    linkcolor=KVGrau, % einfache interne Verknüpfungen
    anchorcolor=KVGrau,% Ankertext
    citecolor=KVGrau, % Verweise auf Literaturverzeichniseinträge im Text
    filecolor=KVGrau, % Verknüpfungen, die lokale Dateien öffnen
    menucolor=KVGrau, % Acrobat-Menüpunkte
    urlcolor=KVGrau,
% diese Farbdefinitionen sollten für den Druck verwendet werden (alles schwarz)
    %linkcolor=black, % einfache interne Verknüpfungen
    %anchorcolor=black, % Ankertext
    %citecolor=black, % Verweise auf Literaturverzeichniseinträge im Text
    %filecolor=black, % Verknüpfungen, die lokale Dateien öffnen
    %menucolor=black, % Acrobat-Menüpunkte
    %urlcolor=black,
%
    %backref, % Quellen werden zurück auf ihre Zitate verlinkt
    pdftex,
    plainpages=false, % zur korrekten Erstellung der Bookmarks
    pdfpagelabels=true, % zur korrekten Erstellung der Bookmarks
    hypertexnames=false, % zur korrekten Erstellung der Bookmarks
    linktocpage % Seitenzahlen anstatt Text im Inhaltsverzeichnis verlinken
]{hyperref}
% Befehle, die Umlaute ausgeben, führen zu Fehlern, wenn sie hyperref als Optionen übergeben werden
\hypersetup{
    pdftitle={\titel -- \untertitel},
    pdfauthor={\autorName},
    pdfcreator={\autorName},
    pdfsubject={\titel -- \untertitel},
    pdfkeywords={\titel -- \untertitel},
}


% zum Einbinden von Programmcode -----------------------------------------------
\usepackage{listings}
\usepackage{xcolor}
\definecolor{hellgelb}{rgb}{1,1,0.9}
\definecolor{colKeys}{rgb}{0,0,1}
\definecolor{colIdentifier}{rgb}{0,0,0}
\definecolor{colComments}{rgb}{0,0.5,0}
\definecolor{colString}{rgb}{1,0,0}
\lstset{
    float=hbp,
	basicstyle=\footnotesize,
    identifierstyle=\color{colIdentifier},
    keywordstyle=\color{colKeys},
    stringstyle=\color{colString},
    commentstyle=\color{colComments},
    backgroundcolor=\color{hellgelb},
    columns=flexible,
    tabsize=2,
    frame=single,
    extendedchars=true,
    showspaces=false,
    showstringspaces=false,
    numbers=left,
    numberstyle=\tiny,
    breaklines=true,
    breakautoindent=true,
	captionpos=b,
}
\lstdefinelanguage{cs}{
	sensitive=false,
	morecomment=[l]{//},
	morecomment=[s]{/*}{*/},
	morestring=[b]",
	morekeywords={
		abstract,event,new,struct,as,explicit,null,switch
		base,extern,object,this,bool,false,operator,throw,
		break,finally,out,true,byte,fixed,override,try,
		case,float,params,typeof,catch,for,private,uint,
		char,foreach,protected,ulong,checked,goto,public,unchecked,
		class,if,readonly,unsafe,const,implicit,ref,ushort,
		continue,in,return,using,decimal,int,sbyte,virtual,
		default,interface,sealed,volatile,delegate,internal,short,void,
		do,is,sizeof,while,double,lock,stackalloc,
		else,long,static,enum,namespace,string},
}
\lstdefinelanguage{natural}{
	sensitive=false,
	morecomment=[l]{/*},
	morestring=[b]",
	morestring=[b]',
	alsodigit={-,*},
	morekeywords={
		DEFINE,DATA,LOCAL,END-DEFINE,WRITE,CALLNAT,PARAMETER,USING,
		IF,NOT,END-IF,ON,*ERROR-NR,ERROR,END-ERROR,ESCAPE,ROUTINE,
		PERFORM,SUBROUTINE,END-SUBROUTINE,CONST,END-FOR,END,FOR,RESIZE,
		ARRAY,TO,BY,VALUE,RESET,COMPRESS,INTO,EQ},
}
\lstdefinelanguage{php}{
	sensitive=false,
	morecomment=[l]{/*},
	morestring=[b]",
	morestring=[b]',
	alsodigit={-,*},
	morekeywords={
		abstract,and,array,as,break,case,catch,cfunction,class,clone,const,
		continue,declare,default,do,else,elseif,enddeclare,endfor,endforeach,
		endif,endswitch,endwhile,extends,final,for,foreach,function,global,
		goto,if,implements,interface,instanceof,namespace,new,old_function,or,
		private,protected,public,static,switch,throw,try,use,var,while,xor
		die,echo,empty,exit,eval,include,include_once,isset,list,require,
		require_once,return,print,unset},
}
 % verwendete Packages
% !TEX root = ../Projektdokumentation.tex

% Seitenränder -----------------------------------------------------------------
\setlength{\topskip}{\ht\strutbox} % behebt Warnung von geometry
\geometry{a4paper,left=25mm,right=15mm,top=25mm,bottom=35mm}

\usepackage[
	automark, % Kapitelangaben in Kopfzeile automatisch erstellen
	headsepline, % Trennlinie unter Kopfzeile
	ilines % Trennlinie linksbündig ausrichten
]{scrpage2}

% Kopf- und Fußzeilen ----------------------------------------------------------
\pagestyle{scrheadings}
% chapterpagestyle gibt es nicht in scrartcl
%\renewcommand{\chapterpagestyle}{scrheadings}
\clearscrheadfoot

% Kopfzeile
\renewcommand{\headfont}{\normalfont} % Schriftform der Kopfzeile
\ihead{\large{\textsc{\titel}}\\ \small{\untertitel} \\[2ex] \textit{\headmark}}
\chead{}
\ohead{\includegraphics[scale=0.2]{\betriebLogo}}
\setlength{\headheight}{15mm} % Höhe der Kopfzeile
%\setheadwidth[0pt]{textwithmarginpar} % Kopfzeile über den Text hinaus verbreitern (falls Logo den Text überdeckt)

% Fußzeile
\ifoot{\autorName}
\cfoot{}
\ofoot{\pagemark}

% Überschriften nach DIN 5008 in einer Fluchtlinie
% ------------------------------------------------------------------------------

% Abstand zwischen Nummerierung und Überschrift definieren
% > Schön wäre hier die dynamische Berechnung des Abstandes in Abhängigkeit
% > der Verschachtelungstiefe des Inhaltsverzeichnisses
\newcommand{\headingSpace}{1.5cm}

% Abschnittsüberschriften im selben Stil wie beim Inhaltsverzeichnis einrücken
\renewcommand*{\othersectionlevelsformat}[3]{
  \makebox[\headingSpace][l]{#3\autodot}
}

% Für die Einrückung wird das Paket tocloft benötigt
%\cftsetindents{chapter}{0.0cm}{\headingSpace}
\cftsetindents{section}{0.0cm}{\headingSpace}
\cftsetindents{subsection}{0.0cm}{\headingSpace}
\cftsetindents{subsubsection}{0.0cm}{\headingSpace}
\cftsetindents{figure}{0.0cm}{\headingSpace}
\cftsetindents{table}{0.0cm}{\headingSpace}


% Allgemeines
% ------------------------------------------------------------------------------

\onehalfspacing % Zeilenabstand 1,5 Zeilen
\frenchspacing % erzeugt ein wenig mehr Platz hinter einem Punkt

% Schusterjungen und Hurenkinder vermeiden
\clubpenalty = 10000
\widowpenalty = 10000
\displaywidowpenalty = 10000

% Quellcode-Ausgabe formatieren
\lstset{numbers=left, numberstyle=\tiny, numbersep=5pt, breaklines=true}
\lstset{emph={square}, emphstyle=\color{red}, emph={[2]root,base}, emphstyle={[2]\color{blue}}}

\counterwithout{footnote}{section} % Fußnoten fortlaufend durchnummerieren
\setcounter{tocdepth}{\subsubsectionlevel} % im Inhaltsverzeichnis werden die Kapitel bis zum Level der subsubsection übernommen
\setcounter{secnumdepth}{\subsubsectionlevel} % Kapitel bis zum Level der subsubsection werden nummeriert

% Aufzählungen anpassen
\renewcommand{\labelenumi}{\arabic{enumi}.}
\renewcommand{\labelenumii}{\arabic{enumi}.\arabic{enumii}.}
\renewcommand{\labelenumiii}{\arabic{enumi}.\arabic{enumii}.\arabic{enumiii}}

% Tabellenfärbung:
\definecolor{heading}{rgb}{0.64,0.78,0.86}
\definecolor{odd}{rgb}{0.9,0.9,0.9}
 % Definitionen zum Aussehen der Seiten
% !TEX root = ../Projektdokumentation.tex

% Abkürzungen, ggfs. mit korrektem Leerraum
\newcommand{\bs}{$\backslash$\xspace}
\newcommand{\bspw}{bspw.\xspace}
\newcommand{\bzw}{bzw.\xspace}
\newcommand{\ca}{ca.\xspace}
\newcommand{\dahe}{\mbox{d.\,h.}\xspace}
\newcommand{\etc}{etc.\xspace}
\newcommand{\eur}[1]{\mbox{#1\,\texteuro}\xspace}
\newcommand{\evtl}{evtl.\xspace}
\newcommand{\ggfs}{ggfs.\xspace}
\newcommand{\Ggfs}{Ggfs.\xspace}
\newcommand{\gqq}[1]{\glqq{}#1\grqq{}}
\newcommand{\inkl}{inkl.\xspace}
\newcommand{\insb}{insb.\xspace}
\newcommand{\ua}{\mbox{u.\,a.}\xspace}
\newcommand{\usw}{usw.\xspace}
\newcommand{\Vgl}{Vgl.\xspace}
\newcommand{\zB}{\mbox{z.\,B.}\xspace}

% Befehle für häufig anfallende Aufgaben
\newcommand{\Abbildung}[1]{\autoref{fig:#1}}
\newcommand{\Anhang}[1]{\appendixname{}~\ref{#1}: \nameref{#1} \vpageref{#1}}
\newcommand{\includegraphicsKeepAspectRatio}[2]{\includegraphics[width=#2\textwidth,height=#2\textheight,keepaspectratio]{#1}}
\newcommand{\Zitat}[2][\empty]{\ifthenelse{\equal{#1}{\empty}}{\citep{#2}}{\citep[#1]{#2}}}
\newcommand{\Autor}[1]{\textsc{#1}} % zum Ausgeben von Autoren
\newcommand{\itemd}[2]{\item{\textbf{#1}}\\{#2}} % erzeugt ein Listenelement mit fetter Überschrift

% fügt Tabellen aus einer TEX-Datei ein
\newcommand{\tabelle}[3] % Parameter: caption, label, file
{\begin{table}[htbp]
\centering
\singlespacing
\input{Tabellen/#3}
\caption{#1}
\label{#2}
\end{table}}

\newcommand{\tabelleAnhang}[1] % Parameter: file
{\begin{center}
\singlespacing
\input{Tabellen/#1}
\end{center}}

% einfaches Wechseln der Schrift, z.B.: \changefont{cmss}{sbc}{n}
\newcommand{\changefont}[3]{\fontfamily{#1} \fontseries{#2} \fontshape{#3} \selectfont}

% Verwendung analog zu \includegraphics
\newlength{\myx} % Variable zum Speichern der Bildbreite
\newlength{\myy} % Variable zum Speichern der Bildhöhe
\newcommand\includegraphicstotab[2][\relax]{%
% Abspeichern der Bildabmessungen
\settowidth{\myx}{\includegraphics[{#1}]{#2}}%
\settoheight{\myy}{\includegraphics[{#1}]{#2}}%
% das eigentliche Einfügen
\parbox[c][1.1\myy][c]{\myx}{%
\includegraphics[{#1}]{#2}}%
}

\definecolor{AOBlau}{rgb}{0, 0.28, 0.56}

% verschiedene Befehle um Wörter semantisch auszuzeichnen ----------------------
\newcommand{\Index}[2][\empty]{\ifthenelse{\equal{#1}{\empty}}{\index{#2}#2}{\index{#1}#2}}
\newcommand{\Fachbegriff}[2][\empty]{\ifthenelse{\equal{#1}{\empty}}{\textit{\Index{#2}}}{\textit{\Index[#1]{#2}}}}
\newcommand{\NeuerBegriff}[2][\empty]{\ifthenelse{\equal{#1}{\empty}}{\textbf{\Index{#2}}}{\textbf{\Index[#1]{#2}}}}

\newcommand{\Ausgabe}[1]{\texttt{#1}}
\newcommand{\Eingabe}[1]{\texttt{#1}}
\newcommand{\Code}[1]{\texttt{#1}}
\newcommand{\Datei}[1]{\texttt{#1}}

\newcommand{\Assembly}[1]{\textsf{#1}}
\newcommand{\Klasse}[1]{\textsf{#1}}
\newcommand{\Methode}[1]{\textsf{#1}}
\newcommand{\Attribut}[1]{\textsf{#1}}

\newcommand{\Datentyp}[1]{\textsf{#1}}
\newcommand{\XMLElement}[1]{\textsf{#1}}
\newcommand{\Webservice}[1]{\textsf{#1}}

\newcommand{\Refactoring}[1]{\Fachbegriff{#1}}
\newcommand{\CodeSmell}[1]{\Fachbegriff{#1}}
\newcommand{\Metrik}[1]{\Fachbegriff{#1}}
\newcommand{\DesignPattern}[1]{\Fachbegriff{#1}}
 % eigene allgemeine Befehle, die z.B. die Arbeit mit LaTeX erleichtern
% Tabellen, die den Zwischenstand nach einer Projektphase verdeutlichen
\newcommand{\Zwischenstand}[2]{\subsubsection{Zwischenstand}Tabelle~\ref{tab:#2} zeigt den Zwischenstand nach der #1.\tabelle{Zwischenstand nach der #1}{tab:#2}{#2.tex}}

% Abkürzungen
\newcommand{\projektName}{\acs{SSRS} Statistics\xspace}
\newcommand{\teamName}{\ac{DWH} \& \ac{VOM}\xspace}
 % eigene projektspezifische Befehle, z.B. Abkürzungen usw.

\begin{document}

% kann nach dem Lesen entfernt werden ---------------------------------------
\pagestyle{plain}
% !TEX root = Projektdokumentation.tex
\section*{Über diese Vorlage}

Diese \LaTeX-Vorlage wurde von Stefan Macke\footnote{Blog des Autors:
\url{http://fachinformatiker-anwendungsentwicklung.net}, Twitter:
\Eingabe{@StefanMacke}} als Grundlage für die Projektdokumentationen der Auszubildenden zum Fachinformatiker mit Fachrichtung
Anwendungsentwicklung bei der \AO entwickelt. Nichtsdestotrotz dürfte sie ebenso für die anderen IT-Berufe\footnote{\zB IT-Kaufleute, Fachinformatiker
mit Fachrichtung Systemintegration \usw} geeignet sein, da diese anhand der gleichen Verordnung bewertet werden.

Diese Vorlage enthält bereits eine Vorstrukturierung der möglichen Inhalte einer tatsächlichen Projektdokumentation, die auf Basis der
Erfahrungen im Rahmen der Prüfertätigkeit des Autors erstellt und unter Zuhilfenahme von \citet{Rohrer2011} abgerundet wurden.

Sämtliche verwendeten Abbildungen, Tabellen und Listings stammen von \citet{Grashorn2010}.

Download-Link für diese Vorlage: \url{http://fiae.link/LaTeXVorlageFIAE}

Auch verfügbar auf GitHub: \url{https://github.com/StefanMacke/latex-vorlage-fiae}

\subsection*{Lizenz}

\begin{center}
\includegraphicsKeepAspectRatio{CC-Logo.pdf}{0.3}
\end{center}
Dieses Werk steht unter einer Creative Commons Namensnennung - Weitergabe unter gleichen Bedingungen 4.0 International Lizenz.
\footnote{\url{http://creativecommons.org/licenses/by-sa/4.0/}}

\begin{center}
\includegraphicsKeepAspectRatio{CC-Attribution.pdf}{0.07}
\includegraphicsKeepAspectRatio{CC-ShareAlike.pdf}{0.07}
\end{center}

\begin{description}
	\item[Namensnennung] Sie müssen den Namen des Autors/Rechteinhabers in der von ihm festgelegten Weise nennen.
	\footnote{Die Namensnennung im \LaTeX-Quelltext mit Link auf \url{http://fiae.link/LaTeXVorlageFIAE} reicht hierfür aus.}
	\item[Weitergabe unter gleichen Bedingungen] Wenn Sie das lizenzierte Werk \bzw den lizenzierten Inhalt bearbeiten
	oder in anderer Weise erkennbar als Grundlage für eigenes Schaffen verwenden, dürfen Sie die daraufhin neu entstandenen
	Werke \bzw Inhalte nur unter Verwendung von Lizenzbedingungen weitergeben, die mit denen dieses Lizenzvertrages identisch oder vergleichbar sind.
\end{description}

\subsection*{Inhalt der Projektdokumentation}

Grundsätzlich definiert die \citet[S.~1746]{Bundesgesetzblatt48}\footnote{Dieses
Dokument sowie alle weiteren hier genannten können unter
\url{http://fiae.link/LaTeXVorlageFIAEQuellen} heruntergeladen werden.} das Ziel der Projektdokumentation wie folgt:
\begin{quote}
"`Durch die Projektarbeit und deren Dokumentation soll der Prüfling belegen, daß er Arbeitsabläufe und Teilaufgaben zielorientiert unter
Beachtung wirtschaftlicher, technischer, organisatorischer und zeitlicher Vorgaben selbständig planen und kundengerecht umsetzen sowie
Dokumentationen kundengerecht anfertigen, zusammenstellen und modifizieren kann."'
\end{quote}

Und das \citet[S.~36]{BMBF2000} ergänzt:
\begin{quote}
"`Die Ausführung der Projektarbeit wird mit praxisbezogenen Unterlagen dokumentiert.
Der Prüfungsausschuss bewertet die Projektarbeit anhand der Dokumentation. Dabei
wird nicht das Ergebnis -- \zB ein lauffähiges Programm -- herangezogen, sondern
der Arbeitsprozess. Die Dokumentation ist keine wissenschaftliche Abhandlung,
sondern eine handlungsorientierte Darstellung des Projektablaufs mit
praxisbezogenen, d.h. betriebüblichen Unterlagen. Sie soll einen Umfang von
maximal 10 bis 15 DIN A 4-Seiten nicht überschreiten. Soweit erforderlich können in
einem Anhang \zB den Zusammenhang erläuternde Darstellungen beigefügt werden."'
\end{quote}

Außerdem werden dort die grundlegenden Inhalte der Projektdokumentation aufgelistet:
\begin{itemize}
	\item Name und Ausbildungsberuf des Prüfungsteilnehmers
	\item Angabe des Ausbildungsbetriebes
	\item Thema der Projektarbeit
	\item Falls erforderlich, Beschreibung/Konkretisierung des Auftrages
	\item Umfassende Beschreibung der Prozessschritte und der erzielten Ergebnisse
	\item Gegebenenfalls Veränderungen zum Projektantrag mit Begründung
	\item Wenn für das Projekt erforderlich, ein Anhang mit praxisbezogenen Unterlagen und Dokumenten. Dieser Anhang sollte nicht
	aufgebläht werden. Die angehängten Dokumente und Unterlagen sind auf das absolute Minimum zu beschränken.
\end{itemize}

In den folgenden Kapiteln werden diese geforderten Inhalte und sinnvolle Ergänzungen nun meist stichwortartig und \ggfs mit
Beispielen beschrieben. Nicht alle Kapitel müssen in jeder Dokumentation vorhanden sein. Handelt es sich \bspw um ein in sich
geschlossenes Projekt, kann das Kapitel~\ref{sec:Projektabgrenzung}: \nameref{sec:Projektabgrenzung} entfallen; arbeitet die
Anwendung nur mit \acs{XML}-Dateien, kann und muss keine Datenbank beschrieben werden \usw


\subsection*{Formale Vorgaben}

Die formalen Vorgaben zum Umfang und zur Gestaltung der Projektdokumentation können je nach IHK recht unterschiedlich sein.
Normalerweise sollte die zuständige IHK einen Leitfaden bereitstellen, in dem alle Formalien nachgelesen werden können,
wie \zB bei der \citet{MerkblattIHK}.

Als Richtwert verwende ich 15 Seiten für den reinen Inhalt. Also in dieser Vorlage alle Seiten, die arabisch nummeriert
sind (ohne das Literaturverzeichnis und die eidesstattliche Erklärung).
Große Abbildungen, Quelltexte, Tabellen \usw gehören in den Anhang, der 25 Seiten nicht überschreiten sollte.

Typographische Konventionen, Seitenränder \usw können in der Datei \Datei{Seitenstil.tex} beliebig angepasst werden.


\subsection*{Bewertungskriterien}
Die Bewertungskriterien für die Benotung der Projektdokumentation sind recht einheitlich und können leicht in Erfahrung
gebracht werden, \zB bei der \citet{BewertungsmatrikIHK}.
Grundsätzlich sollte die Projektdokumentation nach der Fertigstellung noch einmal im Hinblick auf diese Kriterien durchgeschaut werden.

\cleardoublepage
\pagestyle{scrheadings}
% ---------------------------------------------------------------------------

\phantomsection
\thispagestyle{empty}
\pdfbookmark[1]{Eidesstattliche Erklärung}{ihkdeckblatt}
\includegraphicsKeepAspectRatio{DeckblattIHK}{1}
\cleardoublepage

\phantomsection
\thispagestyle{plain}
\pdfbookmark[1]{Deckblatt}{deckblatt}
% !TEX root = Projektdokumentation.tex
\begin{titlepage}

\begin{center}
\includegraphics[scale=0.2]{LogoIHK.pdf}\\[1ex]
\Large{Abschlussprüfung \pruefungstermin}\\[3ex]

\Large{\ausbildungsberuf}\\
\LARGE{\betreff}\\[4ex]

\huge{\textbf{\titel}}\\[1.5ex]
\Large{\textbf{\untertitel}}\\[2ex]

\normalsize
Abgabetermin: \abgabeOrt, den \abgabeTermin\\[1em]
\textbf{Prüfungsbewerber:}\\
\autorName\\
\autorAnschrift\\
\autorOrt\\[5ex]

\includegraphics[scale=0.3]{\betriebLogo}\\[2ex]
\textbf{Ausbildungsbetrieb:}\\
\betriebName\\
\betriebAnschrift\\
\betriebOrt\\[3em]
\end{center}

\small
\noindent
Dieses Werk einschließlich seiner Teile ist \textbf{urheberrechtlich geschützt}.
Jede Verwertung außerhalb der engen Grenzen des Urheberrechtgesetzes ist ohne
Zustimmung des Autors unzulässig und strafbar. Das gilt insbesondere für
Vervielfältigungen, Übersetzungen, Mikroverfilmungen sowie die Einspeicherung
und Verarbeitung in elektronischen Systemen.

\end{titlepage}
\cleardoublepage

% Preface --------------------------------------------------------------------
\phantomsection
\pagenumbering{Roman}
\pdfbookmark[1]{Inhaltsverzeichnis}{inhalt}
\tableofcontents
\cleardoublepage

\phantomsection
\listoffigures
\cleardoublepage

\phantomsection
\listoftables
\cleardoublepage

\phantomsection
\lstlistoflistings
\cleardoublepage

\newcommand{\abkvz}{Abkürzungsverzeichnis}
\renewcommand{\nomname}{\abkvz}
\section*{\abkvz}
\markboth{\abkvz}{\abkvz}
\addcontentsline{toc}{section}{\abkvz}
% !TEX root = Projektdokumentation.tex

% Es werden nur die Abkürzungen aufgelistet, die mit \ac definiert und auch benutzt wurden. 
%
% \acro{VERSIS}{Versicherungsinformationssystem\acroextra{ (Bestandsführungssystem)}}
% Ergibt in der Liste: VERSIS Versicherungsinformationssystem (Bestandsführungssystem)
% Im Text aber: \ac{VERSIS} -> Versicherungsinformationssystem (VERSIS)

% Hinweis: allgemein bekannte Abkürzungen wie z.B. bzw. u.a. müssen nicht ins Abkürzungsverzeichnis aufgenommen werden
% Hinweis: allgemein bekannte IT-Begriffe wie Datenbank oder Programmiersprache müssen nicht erläutert werden,
%          aber ggfs. Fachbegriffe aus der Domäne des Prüflings (z.B. Versicherung)

% Die Option (in den eckigen Klammern) enthält das längste Label oder
% einen Platzhalter der die Breite der linken Spalte bestimmt.
\begin{acronym}[WWWWW]
	\acro{AO}{\textsc{Alte Oldenburger} Krankenversicherung AG}
	\acro{API}{Application Programming Interface}
	\acro{CSS}{Cascading Style Sheets}
	\acro{CSV}{Comma Separated Value}
	\acro{EPK}{Ereignisgesteuerte Prozesskette}
	\acro{ERM}{En\-ti\-ty-Re\-la\-tion\-ship-Mo\-dell}
	\acro{HTML}{Hypertext Markup Language}\acused{HTML}
	\acro{IDE}{Integrated Development Environment}
	\acro{MVC}[MVC]{Model View Controller}
	\acro{NatInfo}[\textsc{NatInfo}]{Natural Information System}
	\acro{Natural}[\textsc{Natural}]{Programmiersprache der Software AG}\acused{Natural}
	\acro{ORM}{Object-Relational Mapping}
	\acro{PHP}{Hypertext Preprocessor}
	\acro{SDK}{Software Development Kit}
	\acro{SQL}{Structured Query Language}
	\acro{SVN}{Subversion}
	\acro{UML}{Unified Modeling Language}
	\acro{XML}{Extensible Markup Language}
\end{acronym}

\clearpage

% Inhalt ---------------------------------------------------------------------
\pagenumbering{arabic}
% !TEX root = Projektdokumentation.tex
% !TEX root = ../Projektdokumentation.tex
\section{Einleitung}
\label{sec:Einleitung}
Diese Projektdokumentation zu der Webanwendung \projektName für die \ac{KVWL}, welche im Rahmen der IHK-Abschlussprüfung Sommer 2019 erstellt wurde, beschreibt umfassend \ua die Thematik, notwendige Prozessschritte und erzielte Ergebnisse des Projekts. Zum besseren Verständnis erhält sie sowohl ergänzende, praxisbezogene Unterlagen und Dokumente (wie Tabellen, Grafiken, Screenshots und Quellcode) als auch eine Kundendokumentation, die Bestandteil des Anhangs ist.

\subsection{Projektumfeld} 
\label{sec:Projektumfeld}
Die Webapplikation wurde im Rahmen eines internen Projekts bei der \ac{KVWL} in Dortmund umgesetzt. Sie wurde während des Projektzeitraums vollständig auf Basis der Wünsche des Auftraggebers neu entwickelt und soll in dessen Umgebung/Umfeld nach Projektende produktiv eingesetzt werden.

\subsubsection{Unternehmen} 
\label{sec:Unternehmen}
Laut offizieller Unternehmenswebsite\footnote{\cite{KVWL}} vertritt die \ac{KVWL}
\begin{displayquote}
	die Interessen von über 15.000 niedergelassenen Vertragsärzten und Ärztlichen und Psychologischen Psychotherapeuten sowie Kinder- und Jugendlichenpsychotherapeuten in Westfalen-Lippe. Mit rund 2000 Mitarbeitern, von denen ca. die Hälfte im Notfalldienst überwiegend geringfügig beschäftigt ist, ist die \ac{KVWL} eine der größten Kassenärztlichen Vereinigungen. Alle zur vertragsärztlichen Versorgung zugelassenen Ärzte, Ärztlichen und Psychologischen Psychotherapeuten, sowie Kinder- und Jugendlichenpsychotherapeuten sind Pflichtmitglieder der \ac{KVWL}. Zu den zentralen Aufgaben der \ac{KVWL} zählt die Sicherstellung der ambulanten ärztlichen Versorgung. [...]\\
	Als öffentlich-rechtliche Körperschaft handelt die Kassenärztliche Vereinigung in Vertretung ihrer Mitglieder mit den Verbänden der Krankenkassen die Gesamtvergütung für die ärztlichen Leistungen aus. Dazu schließt sie entsprechende Verträge ab. Als genossenschaftlicher Zusammenschluss garantiert die \ac{KVWL} somit, dass die wirtschaftlichen Interessen ihrer Mitglieder gegenüber den Krankenkassen gewahrt werden.
\end{displayquote}
Darüber hinaus trägt sie dafür Sorge, dass eine ausreichende ärztliche Versorgung sichergestellt ist (Sicherstellungsauftrag) und die Abrechnungen ihrer Mitglieder sachlich-rechnerisch korrekt sind (Gewährleistungsauftrag), wofür unterstützend auch eine Stelle zur Bekämpfung rechtwidriger Finanzmittelnutzung im Gesundheitswesen eingerichtet wurde.\footnote{\Vgl \cite{KVWL}}

\subsubsection{Auftraggeber} 
\label{sec:Auftraggeber}
Wie in dem Projektantrag\footnote{\cite{Projektantrag}} bereits beschrieben ist der
\begin{displayquote}
	Auftraggeber für dieses Projekt [...] das \acs{CFT} \teamName des Geschäftsbereich Informationstechnologie der \ac{KVWL} (\acs{CFT} steht für \aclu{CFT}).
	[...]
	Das IT-Team \teamName verwaltet die Daten der \ac{KVWL} mithilfe unterschiedlicher Datenbankmanagementsysteme wie SQL Server oder Oracle Datenbanken.
\end{displayquote}
Der Begriff \ac{DWH} beschreibt dabei eine zentrale Datenbank, welche Daten aus mehreren heterogenen Quellen zur Zwecken der Analyse zusammenführt und verdichtet.\footnote{\Vgl \cite{wiki:dwh}}
Folglich hat das Team die die Zielsetzung alle im Unternehmen produzierten Daten für Auswertungszwecke an einem zentralen Ort zu verwalten und in einer Form aufzuarbeiten, sodass einem möglichst großer Benutzerpool einfache Zugriffsmöglichkeiten zur Verfügung gestellt werden.
Zum Tagesgeschäft des \teamName gehört die Erstellung, Anpassung und Administration von Berichten, den sogenannten Reports, die den Fachabteilungen für operative und strategische Auswertungen und Analysen bereitgestellt werden.

\subsection{Ausgangssituation} 
\label{sec:Ausgangssituation}
Dieser Abschnitt liefert nur eine kurzen Überblick über die organisatorische und technische Ausgangssituation vor Projektbeginn. Genaueres wird später unter dem Punkt \ref{sec:Anforderungsanalyse} (\ua die Ist- und Soll-Zustände) erläutert werden.

\subsubsection{Organisatorisch}
\label{sec:Ausgangssituation:Organisatorisch}
Da eine der Hauptaufgaben des Teams die Verwaltung der oben erwähnten Reports darstellt, werden kontinuierlich Berichte zur Förderung der \ac{BI} (deutsch: Geschäftsanalytik) erstellt und modifiziert, um Nutzern und Entwicklern einen komfortablen Zugriff zur einfachen Gewinnung verwertbarer Erkenntnisse zu ermöglichen. Hierzu werden \ac{BI}-Reports mit Tabellen und Diagrammen entwickelt, die Benutzern einen zentralisierten und einheitlichen Zugang zu der Vielzahl der komplexen Datenstrukturen gewähren.\footnote{\Vgl \cite{Projektantrag}}
Um über die Gesamtmenge einen Überblick zu bewahren, existieren Metadaten-Berichte (also Berichte über Berichte), die einige Informationen \zB über die Anzahl und den Status bereitstellen.
Die Anzahl der angelegten Metadaten-Berichte ist jedoch sehr gering und die Informationsdichte nicht ausreichend, aus welchem Grund dieses Projekt genau an dieser Stelle ansetzt.

\subsubsection{Technisch}
\label{sec:Ausgangssituation:Technisch}
Zur Berichtserstellung wird die Berichtserstellungsplattform \ac{SSRS} eingesetzt, die Werkzeuge bietet (wie den Berichts-Generator/Designer), um aussagekräftige und anschauliche Berichte auf einem Webportal veröffentlichen zu können.\footnote{\Vgl \cite{Projektantrag}}
Es ist ein serverbasiertes Generierungssystem für Berichte des Herstellers Microsoft und ist nur ein Tool von vielen aus der Reihe der \ac{BI}-Werkzeuge (\zB auch Crystal Reports).



\subsection[Projektziel]{Projektziel\footnote{\cite{Projektantrag}}}
\label{sec:Projektziel}
\begin{displayquote}
	Ziel der Anwendung ist die Verbesserung der Qualität der Gesamtheit der Berichte, um einen fehlerfreien und schnellen Zugriff auf \ac{BI}-Daten zu jeder Zeit sicherzustellen und die Anzahl der ungenutzten Berichte und Datenquellen schrittweise zu reduzieren. Für diesen Prozess soll die Applikation einen komfortablen Einstieg, aber auch Möglichkeiten für eine unkomplizierte Erweiterung bieten, um \ggfs die Erhebung neuer Statistiken zu begünstigen.
\end{displayquote}


\subsection[Projektbegründung]{Projektbegründung\footnote{\cite{Projektantrag}}}
\label{sec:Projektbegruendung}
\begin{displayquote}
	Die stetig wachsende Anzahl von Berichten erschwert den Überblick über den Status (Fehleranfälligkeit), die Performance und die Abhängigkeiten von Berichten sowie die darangekoppelten Datenquellen. Aktuell existieren über 1000 Berichte zur Auswertung diverser Daten für viele Abteilungen, die \ua Transparenz und Zugang zu \ac{KP}-Indikatoren schaffen und so einen erheblichen, produktiven Mehrwert erzeugen. Dennoch fehlen Metadaten über diese Plattform, die Aufwände für die Erstellung, Überarbeitung und Korrektur von Berichten auf die relevanten Stellen fokussieren.
\end{displayquote}

\subsection[Projektbeschreibung]{Projektbeschreibung\footnote{\cite{Projektantrag}}} 
\label{sec:Projektbeschreibung}
\begin{displayquote}
	Aufgabe des Prüflings \autorName\xspace ist daher die Implementierung einer Web-Applikation, die geeignete quantitative Metriken über die Ausführung, den Status, die Performance und die Abhängigkeiten zwischen den Berichten anschaulich visualisiert. Es muss eine Übersicht über die am häufigsten ausgeführten, fehlerhaften, langsamen Berichte und nicht mehr referenzierte Datenquellen angezeigt werden können. Die Nutzung der Applikation sollte ausgewählten Anwendern (Team \teamName) vorbehalten sein, \dahe ein Login mit Anbindung an das unternehmensinterne \ac{AD} muss implementiert werden. Das Backend der Anwendung soll in der Programmiersprache Java umgesetzt, das Frontend sollte als ansprechende HTML-Seite mit CSS gemäß der Cooperate-Design-Richtlinien der \ac{KVWL} entworfen werden.
\end{displayquote}

\subsection{Projektschnittstellen} 
\label{sec:Projektschnittstellen}

\subsubsection{Organisatorisch} 
\label{sec:Projektschnittstellen:Organisatorisch}

\paragraph{Projektverantwortlichkeiten} ~\\
\label{p:Projektverantwortlichkeiten}
Für die Analyse, Umsetzung und Planung des gesamten Projekts ist Yannik Queisler zuständig, der seit dem 11. März 2019 Mitarbeiter des \ac{CFT} \teamName ist. Da dieses Projekt nicht nur ein unternehmens- sondern auch abteilungsinternes Projekt ist, entstammen sämtliche Anforderungen ausschließlich aus diesem Team. Eine detailreichere Aufstellung sämtlicher Beteiligter folgt direkt im Anschluss im Abschnitt \nameref{p:Stakeholder}.

\paragraph{Stakeholder} ~\\
\label{p:Stakeholder}
Die folgende Tabelle zeigt alle beteiligten Stakeholder, \dahe alle am Projektergebnis interessierten Personen (deutsch: Anspruchsgruppe).
\tabelle{Stakeholder}{tab:Stakeholder}{Stakeholder}

\paragraph{Endbenutzer} ~\\
\label{p:Endbenutzer}
Die Anwendung \projektName soll von allen Mitgliedern des \ac{CFT} \teamName genutzt werden können. Entsprechende Zugangsberechtigungen zur Beschränkung auf eben diesen Nutzerkreis sind zu implementieren.

\subsubsection{Technisch} 
\label{sec:Projektschnittstellen:Technisch}



\begin{itemize}
	\item Mit welchen anderen Systemen interagiert die Anwendung (technische Schnittstellen)? 
	\item Wer genehmigt das Projekt \bzw stellt Mittel zur Verfügung? 
	\item Wer sind die Benutzer der Anwendung?
	\item Wem muss das Ergebnis präsentiert werden?
\end{itemize}


\subsection{Projektabgrenzung} 
\label{sec:Projektabgrenzung}
\begin{itemize}
	\item Was ist explizit nicht Teil des Projekts (\insb bei Teilprojekten)?
\end{itemize}

% !TEX root = ../Projektdokumentation.tex
\section{Projektplanung} 
\label{sec:Projektplanung}
Dieser Abschnitt wird sich mit der Ressourcen- und Ablaufplanung von \projektName beschäftigen. Dazu werden die eingesetzten Ressourcen mit den entsprechenden Aufgaben, dem Zeitbedarf und den Kosten angegeben werden. Die Planung der Arbeitsschritte wird in strikter Anlehnung an das verwendete Vorgehensmodell erfolgen und hier ebenfalls textuell und visuell wiedergegeben werden.

\subsection{Ressourcenplanung}
\label{sec:Ressourcenplanung}
Die Planung der Ressourcen wird in Personal- und Sachmittelplanung untergliedert, für welche die Aspekte Kosten- und Zeitbedarf anhand ihrer Aufgaben erläutert werden.

\subsubsection{Personal}
\label{sec:Personal}
Orientiert an der bereits vorgestellten Stakeholder-Tabelle in Abschnitt \ref{p:Stakeholder}, folgt nun die Aufwandsaufstellung in Stunden für die einzelnen Projektrollen sortiert nach Aufwand.
\tabelle{Personalplanung}{tab:Personalplanung}{Personalplanung}

\paragraph{Projektverantwortlicher} ~\\
\label{p:Projektverantwortlicher}
Der Hauptaufwand des Projekts (70 Stunden), der \ua die Planung, Durchführung und Dokumentation umfasst, wird vom Projektverantwortlichen \autorName\xspace erbracht. Zur effizienten Nutzung dieser Arbeitsdauer ist eine kontinuierliche Kommunikation mit den verbleibenden Projektrollen von enormer Wichtigkeit.

\paragraph{Ansprechpartnerin} ~\\
\label{p:Ansprechpartnerin}
Den zweitgrößten Aufwand (15 Stunden) trägt die Team-Ansprechpartnerin Anne Schlebusch. Sie übernimmt zum einen die Beantwortung fachlicher und zum anderen die Beratung bei technischen Fragen. Der Großteil des angegebenen Aufwands fließt dabei in die Anforderungserhebung und die Definition erwünschter Datenbankabfragen.

\paragraph{Benutzer} ~\\
\label{p:Benutzer}
Darauf folgt die Gruppe der Benutzer, also das \ac{CFT} \teamName, mit einem Aufwand von fünf Stunden. Innerhalb dieses Zeitrahmens können arbeitsteilig Aufgaben wie Rücksprachen bei der Anforderungserhebung und Test des Endprodukts übernommen werden.

\paragraph{Betreuung Auszubildender} ~\\
\label{p:Betreuung}
Da Mareike Schroll als Mitarbeiterin der Personalabteilung lediglich für organisatorische Fragen bezüglich der IHK-Abschlussprüfung des Prüflings zuständig ist, wendet sie den geringsten Aufwand (1 Stunde) für das Projekt auf.

\subsubsection{Sachmittel}
\label{sec:Sachmittel}

\paragraph{Hardware} ~\\
\label{p:Hardware}
Zur Fertigstellung von \projektName wird ein Standard-Entwicklerarbeitsplatz, ein virtueller Windows-Server für die Bereitstellung der Applikation und der Microsoft SQL Server benötigt. Da dem Auszubildenden Yannik Queisler bereits seit Ausbildungsbeginn alle notwendigen Geräte des Entwicklerarbeitsplatzes (wie PC, Telefon, Bildschirme \etc) zugeteilt wurden, entstehen hierfür keine Kosten. Auch die erforderliche Server-Infrastruktur (also sowohl der Windows-Server als auch der SQL Server) sind bereits in dem Geschäftsbereich IT vorhanden. Es fallen also keine weiteren Kosten für Hardwareanschaffungen an.

\paragraph{Software} ~\\
\label{p:Software}
Da der oben erwähnte Entwicklerarbeitsplatz nicht nur Hardwarekomponenten, sondern auch Software umfasst, sind für die Entwicklung ebenfalls keine Kosten zu veranschlagen. Eine \ac{IDE} (IntelliJ IDEA), Dokumentationswerkzeuge (Office, Confluence) und Diagramm-Tools (Gliffy) stehen bereits zur Verfügung.

\subsection{Ablaufplanung}
\label{sec:Ablaufplanung}

\subsubsection{Vorgehensmodell}
\label{sec:Vorgehensmodell}
Aufgrund der starren zeitlichen Vorgabe, der im Allgemeinen recht kurzen Projektdauer und den im Voraus eindeutigen Benutzeranforderungen wurde das Projekt nach dem Wasserfallmodell durchgeführt. Dies bietet eine einfache Planung durch fest vorgegebene Phasen und eine klare Einschätzung der Kosten im Projekt. Die fünf Phasen des Modells werden im Folgenden anhand der \Abbildung{Waterfall} kurz erklärt.\footnote{\cite{wiki:waterfall}} 
\begin{figure}[htb]
	\centering
	\includegraphicsKeepAspectRatio{Wasserfallmodell.png}{0.5}
	\caption{Wasserfallmodell}
	\label{fig:Waterfall}
\end{figure}

\paragraph{Anforderungsanalyse und -spezifikation} ~\\
\label{p:Anforderungsanalyse}
Zu Beginn werden alle Anforderungen des Auftraggebers erhoben, strukturiert und validiert. Das daraus resultierende Ergebnis wird schriftlich dokumentiert.

\paragraph{Systementwurf} ~\\
\label{p:Systementwurf}
Im Anschluss werden die Basiskomponenten des Systems mit ihren Abhängigkeiten identifiziert und meist in Form von UML-Diagrammen festgehalten.

\paragraph{Implementation} ~\\
\label{p:Implementation}
Darauf folgt die Umsetzung des Entwurfs in eine ausführbare Software, die Implementation. Diese Phase kann auch das Schreiben von Modultests beinhalten.

\paragraph{Überprüfung} ~\\
\label{p:Ueberpruefung}
Nach erfolgreicher Codierung folgt die Überprüfung der Anwendung durch Integrations- und Systemtests.

\paragraph{Wartung} ~\\
\label{p:Wartung}
Nach der Auslieferung wird die Software fortwährend über die gesamte Lebensdauer gewartet. Diese Phase ist für \projektName nicht relevant, da sich das Projekt mit der Entwicklung der Anwendung und nicht mit deren Wartung befasst.

\subsubsection{Projektphasen}
\label{sec:Projektphasen}
Die eben aufgezählten Phasen des Wasserfallmodells wurden für das Projekt noch um einige Phasen verfeinert. Die grobe Aufstellung der Phasen zeigt Tabelle \ref{tab:Zeitplanung}. Die neu hinzugekommenen Phasen werden nach der Tabelle kurz beschrieben.
\tabelle{Zeitplanung}{tab:Zeitplanung}{ZeitplanungKurz}\\
Eine detailliertere Zeitplanung findet sich im \Anhang{app:Zeitplanung}.

\paragraph{Wirtschaftlichkeitsanalyse} ~\\
\label{p:Wirtschaftlichkeitsanalyse}
Im Rahmen der Wirtschaftlichkeitsanalyse werden Kosten und Nutzen gegenübergestellt und begründet unter Berücksichtigung der Amortisationsdauer über eine Eigenentwicklung oder einen Einkauf einer Fremdanwendung entschieden.

\paragraph{Einarbeitung in das Projektumfeld} ~\\
\label{p:EinarbeitungProjektumfeld}
Die ReportServer-Datenbank wird gründlich mitsamt aller Datenstrukturen analysiert und relevante Tabellen zur Erhebung der Metadaten identifiziert.

\paragraph{Abnahme} ~\\
\label{p:Abnahme}
Das Team \teamName verwendet die Applikation in der Produktivumgebung und entscheidet über die vollständige Erfüllung der Akzeptanzkriterien.

\paragraph{Erstellen der Dokumentation} ~\\
\label{p:Dokumentation}
Die hier vorliegende Projektdokumentation wird verfasst.

\subsubsection{Bearbeitungszeitraum}
\label{sec:Bearbeitungszeitraum}
\begin{itemize}
	\item \textbf{Beginn:} 16.03.2019
	\item \textbf{Ende:} 20.05.2019
\end{itemize}
Obwohl die Bearbeitungszeit für eine Projektdauer von 70 Stunden recht großzügig ausgelegt ist, sind die für das Projekt zur Verfügung stehenden 8-Stunden-Tage sehr begrenzt. In den Projektzeitraum fallen sowohl die schriftliche IHK-Abschlussprüfung, das Softwarepraktikum und die Pflichtveranstaltungen der \ac{FH} Dortmund, als auch Arbeitstage bei der \ac{KVWL}, an denen aufgrund anderer Projekte nicht an \projektName gearbeitet werden kann. Für die schriftliche Prüfung sind Vorbereitungstage einzuplanen, Softwarepraktikums- und Pflichtveranstaltungstage sind aus dem Zeitplan komplett zu streichen und Arbeitstage bei der \ac{KVWL} entfallen \ggfs teilweise.
Aufgrunddessen weist das Projekt eine erhöhte Sensibilität für unvorhergesehene Ereignisse auf, die den Verlauf erheblich beeinflussen können.
Kalenderauszüge der drei Monate des Projektzeitraums befinden sich im \Anhang{app:Kalender}.

\subsubsection{Meilensteine}
\label{sec:Meilensteine}
Es folgt eine Auflistung der Zwischenziele des Projekts mit ihren geplanten Eintrittsdaten.
\tabelle{Meilensteine}{tab:Meilensteine}{Meilensteine}

\subsubsection{Abweichungen vom Projektantrag}
\label{sec:AbweichungenProjektantrag}
Im Projektantrag wird in der Projektbeschreibung von einer anschaulichen Visualisierung\footnote{\cite{Projektantrag}} der Berichte gesprochen. Da für die Statistiken nur Tabellen eingesetzt werden, ist dieser Aspekt anfechtbar.
Außerdem wurde bei den Projektphasen die Wirtschaftlichkeitsanalyse (siehe Tabelle \ref{tab:Zeitplanung}) in den Zeitplan aufgenommen. Dies reduziert die Pufferzeit auf eine Stunde.

\newpage

% !TEX root = ../Projektdokumentation.tex
\section{Projektdurchführung} 
\label{sec:Projektdurchführung}
Das Kapitel Projektdurchführung befasst sich mit allen acht, in Abschnitt \ref{sec:Projektphasen} erwähnten, Projektphasen. Prozessschritte und Vorgehensweisen werden differenziert dargestellt und jegliche Abweichungen vom Projektantrag werden umfassend begründet werden. Zudem werden Entscheidungen zu auftretenden Anpassungen und daraus resultierende Folgen erläutert werden.


\subsection{Wirtschaftlichkeitsanalyse}
\label{sec:Wirtschaftlichkeitsanalyse}
Eine Wirtschaftlichkeitsanalyse stellt die Kosten und das erwartete Ergebnis eines Projektes gegenüber. Ziel ist die Überprüfung der Rentabilität \bzw Wirtschaftlichkeit. Gründe für erfolgreiche Projekte sollen identifiziert und die Wirkung von Schwächen begrenzt werden.\footnote{\Vgl \cite{finlex:wa}}\\

\subsubsection{\gqq{Make or Buy}-Entscheidung}
\label{sec:MakeOrBuyEntscheidung}
Eine umfangreiche Recherche innerhalb des eigenen Unternehmens und aller verfügbaren Internetquellen ergab, dass schon mehrere Lösungen zur Darstellung von Statistiken für die Berichterstellungsplattform \ac{SSRS} existieren. Ein repräsentatives Beispiel aller online verfügbaren Ansätze für die Erhebung von Metadaten liefert der Lösungsansatz des Autors Jeff Pries. Er stellt eine Reihe von Informationen über Benutzer, Berichtsausführungen und Performance als Report innerhalb der \ac{SSRS} dar.\footnote{\Vgl \cite{pries:usage}} Da aber vom Auftraggeber eine \gqq{Implementierung einer Web-Applikation}\footnote{\cite{Projektantrag}} gefordert wurde, die unabhängig und außerhalb des Report-Servers besteht, entspricht keine der gefundenen Lösungen den Anforderungen. Folglich wurde die Entscheidung gefällt \projektName in Form einer Individualsoftware zu realisieren. Diese Umsetzung wird außerdem nicht an externe Dienstleister vergeben, sondern an die interne IT-Abteilung, weil diese am besten mit den Rahmenbedingungen vertraut ist. So kann gezielt auf die Bedürfnisse des Auftraggebers \teamName eingegangen werden, um einen hohen Individualisierungsgrad zu erzielen.


\subsubsection{Projektkosten}
\label{sec:Projektkosten}
Die Kosten eines Projekts setzen sich aus Personal- und Ressourcenkosten zusammen. Die Kalkulation der Personalkosten erfolgt dabei nicht durch Auflistung der acht Vorgänge und ihrer geschätzten monetären Aufwände, sondern anhand des je Projektrolle aufgewendeten Zeitaufwands. Die angegebenen Stundensätze wurden nicht eigenständig berechnet, sondern bei der Personalabteilung der \ac{KVWL} erfragt. In den Stundensätzen sind sämtliche zusätzlich anfallende Kostenfaktoren (wie Ressourcen- und Fixkosten) enthalten. Da bei \projektName keine Materialkosten anfallen, repräsentieren die Personalkosten die Gesamtprojektkosten. In der untenstehenden Tabelle \ref{tab:Kostenaufstellung} befindet sich die Aufstellung der Kosten.
\tabelle{Kostenaufstellung}{tab:Kostenaufstellung}{Kostenaufstellung.tex}

\subsubsection{Amortisationsdauer}
\label{sec:Amortisationsdauer}
Durch die einfache Fehlerkennung und die schrittweise Reduktion von ungenutzten Berichten und Datenquellen innerhalb der Webanwendung bietet \projektName signifikante Zeiteinsparungen bei der Verwaltung der Berichte.
\begin{itemize}
	\item Welche monetären Vorteile bietet das Projekt (\zB Einsparung von Lizenzkosten, Arbeitszeitersparnis, bessere Usability, Korrektheit)?
	\item Wann hat sich das Projekt amortisiert?
\end{itemize}

\paragraph{Berechnung der Amortisationsdauer} ~\\
Bei einer Zeiteinsparung von 10 Minuten am Tag für jeden der 8 Anwender und 220 Arbeitstagen im Jahr ergibt sich eine gesamte Zeiteinsparung von 
\begin{eqnarray}
8 \cdot 220 \mbox{ Tage/Jahr} \cdot 10 \mbox{ min/Tag} = 17600 \mbox{ min/Jahr} \approx 293 \mbox{ h/Jahr} 
\end{eqnarray}

Dadurch ergibt sich eine jährliche Einsparung von 
\begin{eqnarray}
293 \mbox{h} \cdot \eur{75}{\mbox{/h}} = \eur{21975}
\end{eqnarray}

Die Amortisationszeit beträgt also $\frac{\eur{4350}}{\eur{21975}\mbox{/Jahr}} \approx 0,2 \mbox{ Jahre} \approx 10 \mbox{ Wochen}$.

\subsection{Einarbeitung in das Projektumfeld}
\label{sec:Einarbeitung}
Der Projektverantwortliche \autorName\xspace setzte sich zunächst mit dem Prozess der Berichtserstellung und -verwaltung des \ac{CFT} \teamName auseinander. Hierzu musste sowohl der Umgang mit dem Berichtsdesigner als auch mit dem zugehörigen Webportal erlernt werden. Außerdem mussten insbesondere die von den \ac{SSRS} persistierten Daten analysiert werden, um Statistiken über jene erfassen zu können.

\subsubsection{Berichtserstellung}
\label{sec:Berichtserstellung}
Die Erstellung von Reports erfolgt bei der \ac{KVWL} mithilfe der \ac{IDE} Visual Studio. Es werden zuerst mit einem Assistenten Datenquellen für den Bericht festgelegt und die gewünschten Tabellen angegeben. Daraufhin kann der Bericht beliebig mit einer Reihe von Grafikelementen (Tabellen, Diagramme, Textfelder \etc) gestaltet werden.

\subsubsection{Berichtsverwaltung}
\label{sec:Berichtsverwaltung}
Zur Verwaltung der Berichte wird ein in SharePoint integriertes Webportal eingesetzt, das alle publizierten Reports in einer Ordnerstruktur auflistet. Die Verzeichnishierarchien sind wie im Windows-Explorer durchsuchbar und die Reports mit der Angabe der benötigten Parameter einfach auszuführen.

\subsubsection{Analyse der Datenstrukturen}
\label{sec:Datenstrukturen}
Die von \ac{SSRS} gesammelten Daten werden in relationalen Tabellen auf dem SQL Server abgelegt. Alle für die Erhebung von Statistiken erforderlichen Tabellen werden in dem nachfolgenden relationalen Datenmodell in \Abbildung{Datenmodell} mitsamt ihrer Abhängigkeiten aufgelistet.
\begin{figure}[htb]
	\centering
	\includegraphicsKeepAspectRatio{Datenmodell.png}{1}
	\caption{Relationales Datenmodell}
	\label{fig:Datenmodell}
\end{figure} 

\subsection{Anwendungsfälle}
\label{sec:Anwendungsfaelle}
\begin{itemize}
	\item Welche Anwendungsfälle soll das Projekt abdecken?
	\item Einer oder mehrere interessante (!) Anwendungsfälle könnten exemplarisch durch ein Aktivitätsdiagramm oder eine \ac{EPK} detailliert beschrieben werden. 
\end{itemize}

\paragraph{Beispiel}
Ein Beispiel für ein Use Case-Diagramm findet sich im \Anhang{app:UseCase}.


\subsection{Qualitätsanforderungen}
\label{sec:Qualitaetsanforderungen}
\begin{itemize}
	\item Welche Qualitätsanforderungen werden an die Anwendung gestellt (\zB hinsichtlich Performance, Usability, Effizienz \etc (siehe \citet{ISO9126}))?
\end{itemize}


\subsection{Lastenheft/Fachkonzept}
\label{sec:Lastenheft}
\begin{itemize}
	\item Auszüge aus dem Lastenheft/Fachkonzept, wenn es im Rahmen des Projekts erstellt wurde.
	\item Mögliche Inhalte: Funktionen des Programms (Muss/Soll/Wunsch), User Stories, Benutzerrollen
\end{itemize}

\paragraph{Beispiel}
Ein Beispiel für ein Lastenheft findet sich im \Anhang{app:Lastenheft}. 

\Zwischenstand{Analysephase}{Analyse}


% !TEX root = ../Projektdokumentation.tex
\section{Abnahmephase} 
\label{sec:Abnahmephase}

\begin{itemize}
	\item Welche Tests (\zB Unit-, Integrations-, Systemtests) wurden durchgeführt und welche Ergebnisse haben sie geliefert (\zB Logs von Unit Tests, Testprotokolle der Anwender)?
	\item Wurde die Anwendung offiziell abgenommen?
\end{itemize}

\paragraph{Beispiel}
Ein Auszug eines Unit Tests befindet sich im \Anhang{app:Test}. Dort ist auch der Aufruf des Tests auf der Konsole des Webservers zu sehen.


\Zwischenstand{Abnahmephase}{Abnahme}


% !TEX root = ../Projektdokumentation.tex
\section{Einführungsphase}
\label{sec:Einfuehrungsphase}

\begin{itemize}
	\item Welche Schritte waren zum Deployment der Anwendung nötig und wie wurden sie durchgeführt (automatisiert/manuell)?
	\item Wurden \ggfs Altdaten migriert und wenn ja, wie?
	\item Wurden Benutzerschulungen durchgeführt und wenn ja, Wie wurden sie vorbereitet?
\end{itemize}


\Zwischenstand{Einführungsphase}{Einfuehrung}

% !TEX root = ../Projektdokumentation.tex
\section{Dokumentation}
\label{sec:Dokumentation}

\begin{itemize}
	\item Wie wurde die Anwendung für die Benutzer/Administratoren/Entwickler dokumentiert (\zB Benutzerhandbuch, \acs{API}-Dokumentation)?
	\item Hinweis: Je nach Zielgruppe gelten bestimmte Anforderungen für die Dokumentation (\zB keine IT-Fachbegriffe in einer Anwenderdokumentation verwenden, aber auf jeden Fall in einer Dokumentation für den IT-Bereich).
\end{itemize}

\paragraph{Beispiel}
Ein Ausschnitt aus der erstellten Benutzerdokumentation befindet sich im \Anhang{app:BenutzerDoku}.
Die Entwicklerdokumentation wurde mittels PHPDoc\footnote{Vgl. \cite{phpDoc}} automatisch generiert. Ein beispielhafter Auszug aus der Dokumentation einer Klasse findet sich im \Anhang{app:Doc}. 

\Zwischenstand{Dokumentation}{Dokumentation}

% !TEX root = ../Projektdokumentation.tex
\section{Fazit} 
\label{sec:Fazit}

\subsection{Soll-/Ist-Vergleich}
\label{sec:SollIstVergleich}

\begin{itemize}
	\item Wurde das Projektziel erreicht und wenn nein, warum nicht?
	\item Ist der Auftraggeber mit dem Projektergebnis zufrieden und wenn nein, warum nicht?
	\item Wurde die Projektplanung (Zeit, Kosten, Personal, Sachmittel) eingehalten oder haben sich Abweichungen ergeben und wenn ja, warum?
	\item Hinweis: Die Projektplanung muss nicht strikt eingehalten werden. Vielmehr sind Abweichungen sogar als normal anzusehen. Sie müssen nur vernünftig begründet werden (\zB durch Änderungen an den Anforderungen, unter-/überschätzter Aufwand).
\end{itemize}

\paragraph{Beispiel (verkürzt)}
Wie in Tabelle~\ref{tab:Vergleich} zu erkennen ist, konnte die Zeitplanung bis auf wenige Ausnahmen eingehalten werden.
\tabelle{Soll-/Ist-Vergleich}{tab:Vergleich}{Zeitnachher.tex}


\subsection{Lessons Learned}
\label{sec:LessonsLearned}

\begin{itemize}
	\item Was hat der Prüfling bei der Durchführung des Projekts gelernt (\zB Zeitplanung, Vorteile der eingesetzten Frameworks, Änderungen der Anforderungen)?
\end{itemize}


\subsection{Ausblick}
\label{sec:Ausblick}

\begin{itemize}
	\item Wie wird sich das Projekt in Zukunft weiterentwickeln (\zB geplante Erweiterungen)?
\end{itemize}


% Literatur ------------------------------------------------------------------
\clearpage
\renewcommand{\refname}{Literaturverzeichnis}
\bibliography{Bibliographie}
\bibliographystyle{Allgemein/natdin} % DIN-Stil des Literaturverzeichnisses
% !TEX root = Projektdokumentation.tex
\clearpage
\addsec{Eidesstattliche Erklärung}

% Hinweis: die eidesstattliche Erklärung ist ggfs. an die Richtlinie der IHK anzupassen

Ich, \autorName, versichere hiermit, dass ich meine \textbf{\betreff} mit dem
Thema
\begin{quote}
\textit{\kompletterTitel}
\end{quote}
selbständig verfasst und keine anderen als die angegebenen Quellen und Hilfsmittel benutzt habe,
wobei ich alle wörtlichen und sinngemäßen Zitate als solche gekennzeichnet habe. Die Arbeit
wurde bisher keiner anderen Prüfungsbehörde vorgelegt und auch nicht veröffentlicht.\\[6ex]

\abgabeOrt, den \abgabeTermin


\rule[-0.2cm]{5.5cm}{0.5pt}

\textsc{\autorName}


% Anhang ---------------------------------------------------------------------
\clearpage
\appendix
\pagenumbering{roman}
% !TEX root = Projektdokumentation.tex
\section{Anhang}

\subsection{Detaillierte Zeitplanung}
\label{app:Zeitplanung}

\subsubsection{Legende}
\label{app:Legende}
Rote Tage sind blockiert (\zB durch Fachhochschulveranstaltungen, Urlaub), gelbe Tage sind nur teilweise verfügbar (weniger als vier Stunden) und grüne Tage sind für das Projekt verwendbar. \\

\subsubsection{Monatskalender}
\label{app:Monatskalender}
\begin{figure}[htb]
    \centering
    \includegraphicsKeepAspectRatio{Zeitplanung_Maerz.png}{0.8}
    \caption{März 2019}
\end{figure}
\begin{figure}[htb]
    \centering
    \includegraphicsKeepAspectRatio{Zeitplanung_April.png}{0.8}
    \caption{April 2019}
\end{figure}
\begin{figure}[htb]
    \centering
    \includegraphicsKeepAspectRatio{Zeitplanung_Mai.png}{0.8}
    \caption{Mai 2019}
\end{figure}

\clearpage

\subsection{Mock-Ups der Anwendung}
\label{app:Mockup}
Im Rahmen des Systementwurfs auf Basis der Anforderungsanalyse wurden Mock-Ups erstellt, die in \Abbildung{Hauptseite} bis \Abbildung{Fehlerstatistiken} gezeigt werden.
\begin{figure}[htb]
    \centering
    \includegraphicsKeepAspectRatio{mock_up_1.pdf}{0.8}
    \caption{Mock-Up: Hauptseite}
    \label{fig:Hauptseite}
\end{figure}
\begin{figure}[htb]
    \centering
    \includegraphicsKeepAspectRatio{mock_up_2.pdf}{0.8}
    \caption{Mock-Up: Ausführungsstatistiken}
\end{figure}
\begin{figure}[htb]
    \centering
    \includegraphicsKeepAspectRatio{mock_up_3.pdf}{0.8}
    \caption{Mock-Up: Abhängigkeitsstatistiken}
\end{figure}
\begin{figure}[htb]
    \centering
    \includegraphicsKeepAspectRatio{mock_up_4.pdf}{0.8}
    \caption{Mock-Up: Performancestatistiken}
\end{figure}
\begin{figure}[htb]
    \centering
    \includegraphicsKeepAspectRatio{mock_up_5.pdf}{0.8}
    \caption{Mock-Up: Fehlerstatistiken}
    \label{fig:Fehlerstatistiken}
\end{figure}

\subsection{Benutzerdokumentation}
\label{app:Benutzerdokumentation}

\end{document}
